\section{Fundamentos matemáticos}\label{tema:mat}
\soporte{F3, F4, F6 y F10}

\concepto{1.1 Monomio y polinomio multivariable}
\que Un monomio en $N$ variables es $c\,x_1^{e_1}x_2^{e_2}\cdots x_N^{e_N}$ con coeficiente $c\in\mathbb{R}$ y exponentes $e_j\in\mathbb{N}$. Un polinomio es una suma de $M$ monomios:
\[
f(\mathbf{x})=\sum_{m=1}^{M} c_m \prod_{j=1}^{N} x_j^{e_{mj}} .
\]
Llamamos $k_m$ al número de variables con $e_{mj}>0$ en el monomio $m$ (sus \emph{factores}) y $F=\sum_m k_m$ al total de factores. Un término constante es un monomio con $k_m=0$.

\porque El enunciado pide funciones \emph{arbitrarias compuestas por múltiples monomios y variables}. Los polinomios son exactamente eso y permiten derivar de forma exacta sin cálculo simbólico general. En casos realistas $k_m\ll N$: cada monomio usa pocas variables. Esa dispersión decide la estructura de datos (tema \ref{tema:datos}) y el coste ($O(F)$ y no $O(MN)$).

\como En el ejemplo $f=3xy+4yz+2x+5$: $N=3$, $M=4$, $k=(2,2,1,0)$ y $F=5$.

\alt Funciones generales (senos, exponenciales) exigirían un árbol de expresión y diferenciación automática o simbólica. Son más generales, pero más lentas y fuera del enunciado, que fija monomios.

\concepto{1.2 Derivada parcial y gradiente}
\que La derivada parcial $\partial f/\partial x_j$ mide la variación de $f$ al mover solo $x_j$. El gradiente es el vector de las $N$ derivadas parciales: $\nabla f=(\partial f/\partial x_1,\dots,\partial f/\partial x_N)$.

\porque Es lo que hay que calcular. Por linealidad de la derivada, el gradiente es la suma de los gradientes de cada monomio. Eso permite repartir monomios entre hilos y sumar contribuciones.

\como Para un monomio, por las reglas del producto y de la potencia:
\[
\frac{\partial}{\partial x_j}\Big(c_m\prod_{l} x_l^{e_{ml}}\Big)=
\begin{cases}
c_m\, e_{mj}\, x_j^{e_{mj}-1}\prod_{l\neq j} x_l^{e_{ml}} & \text{si } e_{mj}>0,\\
0 & \text{si } e_{mj}=0.
\end{cases}
\]
Así, $\partial f/\partial x_j=\sum_{m:\,e_{mj}>0}(\dots)$. En el ejemplo, $\nabla f=(3y+2,\;3x+4z,\;4y)$ y $\nabla f(1,2,3)=(8,15,8)$, que es el caso de prueba obligatorio.

\alt Ninguna para la definición; las alternativas están en cómo calcularla (concepto 1.3).

\concepto{1.3 Métodos para calcular derivadas}
\que Hay cuatro familias:
\begin{itemize}
\item \textbf{Simbólica}: manipula la expresión y obtiene la fórmula (SymPy, Mathematica).
\item \textbf{Analítica estructurada}: aplica la fórmula del concepto 1.2 sobre una estructura conocida (la nuestra).
\item \textbf{Diferencias finitas}: aproxima $\partial f/\partial x_j\approx\big(f(\mathbf{x}+h\mathbf{e}_j)-f(\mathbf{x}-h\mathbf{e}_j)\big)/2h$.
\item \textbf{Diferenciación automática (AD)}: aplica la regla de la cadena a cada operación del programa, en modo directo o inverso.
\end{itemize}

\porque Usamos la analítica estructurada porque es exacta (salvo redondeo), cuesta $O(F)$ y su patrón de acceso es regular, lo que la hace paralelizable. El enunciado pide además automatizar el procedimiento paso a paso del ejemplo, que es justo derivar término a término.

\como El algoritmo concreto está en el tema \ref{tema:sec}. El modo inverso de AD aplicado a un producto $\prod_t p_t$ da los productos prefijo y sufijo que usaremos (ejemplo de Speelpenning, en Griewank y Walther, 2008).

\alt
\begin{itemize}
\item Diferencias finitas: triviales de programar, pero cuestan $2N$ evaluaciones completas ($O(NF)$) y tienen error de truncamiento $O(h^2)$ más error de redondeo $O(\varepsilon/h)$. El $h$ óptimo es del orden de $\varepsilon^{1/3}\approx 6\cdot10^{-6}$ y deja unos 10 dígitos correctos. Sirven como comprobación cruzada, no como base.
\item Simbólica: exacta en aritmética racional, pero muy lenta y con crecimiento de expresiones. La usamos solo como oráculo para casos pequeños.
\item AD con librerías (JAX, PyTorch, Enzyme): general, pero añade dependencias, no se adapta a la estructura dispersa y no es lo que se evalúa.
\item Paso complejo $\operatorname{Im} f(\mathbf{x}+ih\mathbf{e}_j)/h$: exacto a precisión de máquina, pero sigue costando $N$ evaluaciones.
\end{itemize}

\concepto{1.4 Aritmética en coma flotante}
\que IEEE 754 \texttt{double} tiene 53 bits de mantisa: redondeo unitario $u=2^{-53}\approx1.11\cdot10^{-16}$ y rango aproximado $10^{\pm308}$. Cada operación introduce un error relativo de como mucho $u$. La suma \emph{no es asociativa}: $(a+b)+c$ puede diferir de $a+(b+c)$.

\porque Tiene tres consecuencias directas:
\begin{enumerate}
\item Las versiones paralelas suman en otro orden. Su resultado puede diferir del secuencial en los últimos bits, así que la comparación ha de usar una tolerancia, no la igualdad.
\item Los atómicos y la planificación dinámica hacen que el orden cambie entre ejecuciones: el resultado no es determinista.
\item Con exponentes altos se desborda: $1.5^{100}\approx4\cdot10^{17}$ cabe, pero un monomio de 100 factores con $x=2$ y $e=100$ vale $2^{10^4}$, que es \texttt{inf}.
\end{enumerate}

\como La cota de error de una suma de $n$ términos es aproximadamente $(n-1)\,u\sum|t_i|$. Si hay signos mezclados, $\sum|t_i|$ puede ser mucho mayor que $|\sum t_i|$ (cancelación) y el error relativo crece. Por eso el criterio recomendado es
\[
|g_j-g_j^{\text{ref}}|\le \text{rtol}\cdot\max\!\big(1,|g_j^{\text{ref}}|\big),
\]
con $\text{rtol}=10^{-12}$ para \texttt{double} con coeficientes positivos y puntos en $[0.5,1.5]$. Si aparecen cancelaciones, se compara contra $\sum|t_i|$. Para evitar desbordamientos, los generadores deben limitar $e\le5$, $k\le16$ y $x\in[0.5,1.5]$: el peor caso es $1.5^{80}\approx1.2\cdot10^{14}$ y $0.5^{80}\approx8\cdot10^{-25}$, ambos representables.

\begin{trampa}
Comparar con \texttt{==} entre la versión secuencial y la paralela, o usar \texttt{-ffast-math}, que permite reordenar operaciones y rompe la reproducibilidad. Si usas \texttt{-ffast-math}, mide el error antes de adoptarlo.
\end{trampa}

\alt Suma de Kahan o suma por parejas (\emph{pairwise}): reducen el error a casi $O(u)$ a cambio de 2 a 4 veces más operaciones. Es razonable solo si se observan cancelaciones. \texttt{long double} o aritmética racional son útiles solo para el oráculo.

\section{Datos, formatos y generadores}\label{tema:datos}
\soporte{F1, F2 y F3}

\concepto{2.1 Formato de entrada en texto}
\que Un formato textual que reproduce la notación del enunciado: \texttt{f(x, y, z) = 3xy + 4yz + 2x + 5} y ecuaciones como \texttt{3x\^{}2 + yz + z + 5 = 2}.

\porque El enunciado trabaja así y facilita depurar a mano. Para que sea inequívoco con muchas variables, fijamos esta regla: \textbf{una variable es una letra seguida opcionalmente de dígitos} (\texttt{x}, \texttt{y}, \texttt{x12}). Así \texttt{xy} es $x\cdot y$ y \texttt{x12x3\^{}2} es $x_{12}x_3^2$. Un nombre nunca empieza por dígito, de modo que no se confunde con el coeficiente.

\como Gramática (EBNF) que implementa \fichero{poli.c}:
\begin{lstlisting}
ecuacion := expr [ "=" numero ]
expr     := [ "+" | "-" ] termino { ( "+" | "-" ) termino }
termino  := [ numero ] { [ "*" ] variable [ "^" entero ] }
variable := letra { digito }
\end{lstlisting}
Fichero de función: línea 1 con los nombres de las variables (fija el orden del gradiente), línea 2 con la expresión, línea 3 con el punto. Fichero de sistema: línea 1 con las variables, línea 2 con $K$ y $K$ líneas de ecuación; el punto, al final. Si una variable se repite dentro de un monomio (\texttt{xx}), se suman sus exponentes; los monomios repetidos (\texttt{2x + ... + 2x}) se dejan separados porque la derivada es lineal.

\alt
\begin{itemize}
\item CSV de dos tablas (el \fichero{generador.py} de la rama): fácil de generar con Polars, pero más lento de leer en C que un binario y sin la notación del enunciado.
\item Binario (cabecera con $N,M,F$ y los arrays volcados con \texttt{fwrite}): carga casi a la velocidad del disco. Recomendado para los casos grandes de benchmark, siempre que el texto siga siendo el formato de validación.
\item Formato con separador explícito (\texttt{3*x0*x1\^{}2}): más robusto, pero se aleja del enunciado. El parser lo acepta igualmente, porque \texttt{*} es opcional.
\end{itemize}

\concepto{2.2 Parser}
\que El código que convierte el texto en la estructura de datos.

\porque Es la pieza que hace la solución \emph{dinámica}: el número de variables y monomios solo se conoce al leer.

\como Analizador de una pasada con un puntero que avanza (\emph{scanner} manual). Coste $O(L)$ en la longitud del texto, más la búsqueda del nombre. En el prototipo la búsqueda es lineal ($O(N)$ por factor): vale para validar, pero para $N=10^4$ hay que usar una tabla hash o, mejor, extraer el índice del propio nombre (\texttt{x12} es la variable 12). Fragmento central:
\begin{lstlisting}[language=C]
while (*c == '*' || isalpha((unsigned char)*c)) {
    if (*c == '*') { c++; continue; }
    const char *s = c++;
    while (isdigit((unsigned char)*c)) c++;
    int v = buscaVar(s, (int)(c - s), vNombres, nVars);
    if (v < 0) return -1;
    int e = 1;
    if (*c == '^') { c++; e = (int)strtol(c, (char **)&c, 10); }
    ...
}
\end{lstlisting}

\alt \texttt{flex}/\texttt{bison} (potente, pero excesivo para esta gramática), expresiones regulares (\texttt{regex.h}, lentas y frágiles con signos) o parsear en Python y pasar binario a C (rápido de escribir, pero el enunciado pide una función C que represente la función).

\begin{trampa}
Medir el tiempo del parser junto al del gradiente. En casos grandes el parseo de texto puede costar más que el cálculo y ocultaría el speedup. Mide y reporta las dos cosas por separado.
\end{trampa}

\concepto{2.3 Representación densa frente a dispersa (CSR)}
\que
\begin{itemize}
\item \textbf{Densa}: matriz $M\times N$ de exponentes (la decisión inicial del 18-oct).
\item \textbf{COO}: lista de tripletas (monomio, variable, exponente).
\item \textbf{CSR} (\emph{compressed sparse row}): \texttt{vOff[M+1]} indica dónde empiezan los factores de cada monomio en \texttt{vIdVar[F]} y \texttt{vExp[F]}.
\item \textbf{ELL}: cada monomio con $k_{\max}$ huecos.
\item \textbf{CSC}: la transpuesta, agrupada por variable.
\end{itemize}

\porque Usamos CSR porque cuesta $O(M+F)$ en memoria y en recorrido, frente a $O(MN)$ de la densa. Con los tamaños de benchmark propuestos ($k$ uniforme en $[1,16]$, media 8.5):

\begin{center}\small
\begin{tabular}{@{}lrrrrrr@{}}
\toprule
Caso & $N$ & $M$ & $F$ & CSR & Extra dos fases & Densa (\texttt{int32}) \\
\midrule
S & $10^3$ & $10^4$ & $8.5\cdot10^4$ & 0.8 MB & 1.0 MB & 0.04 GB \\
M & $10^4$ & $10^6$ & $8.5\cdot10^6$ & 80 MB & 102 MB & 40 GB \\
L & $10^5$ & $10^7$ & $8.5\cdot10^7$ & 800 MB & 1020 MB & 4000 GB \\
\bottomrule
\end{tabular}
\end{center}

CSR $=8M+4(M+1)+8F$ bytes; la columna \emph{extra dos fases} suma \texttt{vContrib} ($8F$), la CSC ($4F$) y sus desplazamientos ($4(N+1)$). La densa es inviable desde el caso M, y en una GPU de 6 GB solo caben S y M en forma densa (y M ni siquiera). En CSR, el caso L entero ocupa unos 1.8 GB: cabe en 6 GB.

\como Estructura usada en toda la guía (\fichero{poli.h}):
\begin{lstlisting}[language=C]
typedef struct {
    int nVars;
    int nMon;
    int nFact;
    double *vCoef;
    int *vOff;
    int *vIdVar;
    int *vExp;
} Poli;
\end{lstlisting}
El ejemplo $3xy+4yz+2x+5$ queda como \texttt{vCoef = \{3,4,2,5\}}, \texttt{vOff = \{0,2,4,5,5\}}, \texttt{vIdVar = \{0,1,1,2,0\}} y \texttt{vExp = \{1,1,1,1,1\}}.

\alt
\begin{itemize}
\item Densa: acceso trivial y buena si $k_m\approx N$ (funciones muy densas y pequeñas, como el ejemplo). Medir la densidad $F/(MN)$: si supera un 30\,\% aproximadamente, la densa puede competir \pv.
\item ELL: accesos perfectamente alineados en GPU (coalescing), pero desperdicia $M(k_{\max}-\bar k)$ huecos: casi el doble con $k\in[1,16]$.
\item COO: más simple de construir, pero repite el índice de monomio ($4F$ bytes más) y obliga a atómicos.
\item Tipos más pequeños (\texttt{uint8\_t} para exponentes si $e\le255$, \texttt{uint16\_t} para variables si $N\le65535$) reducen tráfico de memoria. Es una optimización válida en un problema limitado por memoria (tema \ref{tema:rend}).
\end{itemize}

\concepto{2.4 AoS frente a SoA}
\que \emph{Array of Structures} (un array de \texttt{struct \{int var; int exp;\}}) frente a \emph{Structure of Arrays} (\texttt{vIdVar[]} y \texttt{vExp[]} separados).

\porque En GPU, hilos consecutivos leen posiciones consecutivas de cada array (coalescing). En CPU, SoA facilita la vectorización SIMD. Usamos SoA.

\como Ya está así en \texttt{Poli}.

\alt AoS de 8 bytes (\texttt{int2}) también se lee bien en GPU y mejora la localidad en CPU. Es una variante a medir en F9, no una decisión crítica.

\concepto{2.5 Generadores de datos}
\que Programas Python que crean funciones y sistemas aleatorios con parámetros $N$, $M$, $k_{\max}$, $e_{\max}$, $K$ y semilla.

\porque Hacen falta casos grandes para medir y casos pequeños con respuesta conocida para validar. El enunciado admite \emph{ficheros de entrada elaborados por el alumno}.

\como Reglas para ambos generadores:
\begin{itemize}
\item Generador de números con semilla explícita (\texttt{numpy.random.default\_rng(semilla)} o \texttt{random.Random(semilla)}), nunca el estado global.
\item Rangos seguros: $c\in[1,9]$, $e\in[1,5]$, $k\in[1,16]$, punto en $[0.5,1.5]$ (concepto 1.4).
\item Variables \texttt{x0} a \texttt{x\{N-1\}}, sin repetir dentro de un monomio (\texttt{rng.sample}).
\item Escritura en texto (validación) y en binario (benchmarks grandes).
\item Generar por bloques, como ya hace \fichero{generador.py}, para no agotar memoria con $M=10^7$.
\end{itemize}

\alt Generar en C (más rápido, pero duplicas lógica). Usar el \fichero{generador.py} existente cambiando \texttt{eMax} a 5, \texttt{cMax} a 16 y los nombres a \texttt{x<id>}, y añadiéndole el punto: es la opción más barata.

\concepto{2.6 Oráculo con SymPy}
\que SymPy es una librería de cálculo simbólico en Python. Deriva la expresión y la evalúa con aritmética racional exacta si el punto es entero o racional.

\porque Da una referencia independiente de nuestro código. Es la mejor defensa ante el requisito eliminatorio de CUDA (\emph{si los resultados no coinciden no se sigue corrigiendo}).

\como Prototipo probado en \fichero{docs/guia/codigo/oraculo.py}. Con puntos enteros en $\{-2,-1,1,2,3\}$ el gradiente de referencia es entero, y la versión C dio error 0 en dos casos ($N=50$, $M=200$, semilla 7; $N=300$, $M=2000$, semilla 11):
\begin{lstlisting}[language=Python]
vSim = sp.symbols(vNombres)
sust = dict(zip(vSim, vPunto))
vGrad = [sp.diff(expr, s).subs(sust) for s in vSim]
\end{lstlisting}

\alt
\begin{itemize}
\item Para casos grandes SymPy es demasiado lento. Se usa un oráculo de dos niveles: el secuencial en C, validado contra SymPy en casos pequeños, pasa a ser la referencia de las versiones paralelas en los grandes.
\item NumPy con \texttt{float128} o \texttt{mpmath} sirven como oráculo de precisión alta.
\item JAX como oráculo de AD: posible, pero añade una dependencia pesada.
\end{itemize}

\section{Algoritmo secuencial}\label{tema:sec}
\soporte{F4 y F5}

\concepto{3.1 Cálculo de potencias enteras}
\que Calcular $x^e$ con $e$ entero.

\porque Es la operación más repetida: dos veces por factor. Su coste domina el cálculo.

\como Exponenciación por cuadrados, $O(\log e)$ multiplicaciones (Knuth, TAOCP vol. 2, 4.6.3). Para la derivada se calcula $b=x^{e-1}$ una sola vez y se obtienen $e\cdot b$ y $x^e=b\cdot x$:
\begin{lstlisting}[language=C]
double powi(double x, int e) {
    double r = 1.0;
    while (e > 0) {
        if (e & 1) r *= x;
        x *= x;
        e >>= 1;
    }
    return r;
}
\end{lstlisting}

\alt
\begin{itemize}
\item \texttt{pow()} de \texttt{math.h}: general (exponentes reales), pero más lenta y con posibles errores de 1 ulp \pv. En CUDA, \texttt{pow} en \texttt{double} es especialmente cara.
\item Bucle de $e$ multiplicaciones: igual de bueno para $e\le3$, peor para exponentes grandes.
\item Tabla precomputada \texttt{mPot[v][e]} para todas las variables y exponentes ($N\cdot e_{\max}$ valores): convierte cada potencia en una lectura. Conviene si $N\cdot e_{\max}\cdot8$ bytes cabe en caché (por ejemplo $10^4\cdot5\cdot8=400$ KB) y si cada variable aparece en muchos monomios. Se mide en F9.
\end{itemize}

\concepto{3.2 Algoritmo del gradiente: ingenuo, división y prefijo-sufijo}
\que Tres formas de obtener, para cada factor $t$ de un monomio, el producto del resto de factores:
\begin{itemize}
\item \textbf{Ingenuo}: recalcula el producto de los otros $k-1$ factores. $O(k^2)$ por monomio.
\item \textbf{División}: calcula $P=\prod_t p_t$ y usa $P/p_t$. $O(k)$, pero falla si algún $x=0$ y pierde precisión cerca de cero.
\item \textbf{Prefijo-sufijo}: $\text{pre}_t=\prod_{s<t}p_s$ y $\text{suf}_t=\prod_{s>t}p_s$, en dos pasadas. $O(k)$, sin divisiones y robusto.
\end{itemize}

\porque Usamos prefijo-sufijo: es el modo inverso de AD aplicado al producto. La diferencia es medible: en la máquina de pruebas el algoritmo ingenuo con 2 hilos fue 3 veces \emph{más lento} que el prefijo-sufijo con 1 hilo (sección \ref{tema:rend}).

\como Implementación de referencia (\fichero{poli.c}, probada con el ejemplo y con el oráculo):
\begin{lstlisting}[language=C]
void gradSec(const Poli *p, const double *vX, double *vGrad) {
    double vPre[KMAX], vDer[KMAX];
    memset(vGrad, 0, sizeof(double) * p->nVars);
    for (int m = 0; m < p->nMon; m++) {
        int ini = p->vOff[m], k = p->vOff[m + 1] - ini;
        double acum = p->vCoef[m];
        for (int t = 0; t < k; t++) {
            double x = vX[p->vIdVar[ini + t]];
            int e = p->vExp[ini + t];
            double base = powi(x, e - 1);
            vPre[t] = acum;
            vDer[t] = e * base;
            acum *= base * x;
        }
        double suf = 1.0;
        for (int t = k - 1; t >= 0; t--) {
            int v = p->vIdVar[ini + t];
            double x = vX[v];
            int e = p->vExp[ini + t];
            vGrad[v] += vPre[t] * vDer[t] * suf;
            suf *= powi(x, e);
        }
    }
}
\end{lstlisting}
Salida del programa \fichero{ejemplo.c}: \texttt{f = 37} y \texttt{grad = (8, 15, 8)}.

\begin{trampa}
\texttt{KMAX} limita los factores por monomio (64 en el prototipo). Si el generador produce $k>$\texttt{KMAX} (el \fichero{generador.py} actual llega a 100), el programa escribe fuera del array. Valida $k_{\max}$ al leer o reserva el buffer dinámicamente.
\end{trampa}

\alt
\begin{itemize}
\item La segunda pasada recalcula \texttt{powi(x, e)}. Guardar también $p_t$ en un tercer buffer ahorra esa potencia a cambio de memoria local. En GPU puede no compensar por la presión de registros: se mide.
\item Horner multivariable reduce multiplicaciones al evaluar $f$, pero requiere reordenar los monomios en un árbol. Es más complejo y paraleliza peor.
\end{itemize}

\concepto{3.3 Coste y memoria del secuencial}
\que Tiempo $O(F\log e_{\max})$ y memoria $O(M+F+N)$.

\porque Permite predecir tiempos y elegir tamaños: el tiempo debe crecer linealmente con $F$. Si no lo hace, hay un problema de caché (por ejemplo, $N$ demasiado grande para que \texttt{vX} quepa en L2).

\como Por factor: 2 llamadas a \texttt{powi} (unas 3 a 6 multiplicaciones cada una con $e\le5$), 4 multiplicaciones y 1 suma; se leen 8 bytes (\texttt{vIdVar} y \texttt{vExp}) más un acceso irregular a \texttt{vX} y otro a \texttt{vGrad}. En la máquina de pruebas, $F=8.5\cdot10^5$ tardó unos 14 ms (unos 16 ns por factor).

\alt No hay algoritmo asintóticamente mejor: cualquier método debe leer los $F$ factores.

\concepto{3.4 C, C++ y Rust}
\que Los tres lenguajes de las versiones secuenciales.

\porque C es la base: OpenMP y CUDA se integran sin intermediarios y das control total de la memoria. C++ aporta \texttt{std::vector} y RAII. Rust aporta seguridad de memoria. Las tres versiones sirven como comprobación cruzada y como comparación de rendimiento.

\como Mantén la misma estructura CSR y el mismo algoritmo en los tres para que la comparación sea justa. Versión Rust probada (\fichero{grad.rs}, salida \texttt{[8.0, 15.0, 8.0]}):
\begin{lstlisting}
fn gradSec(p: &Poli, vX: &[f64]) -> Vec<f64> {
    let mut vGrad = vec![0.0; p.nVars];
    let mut vPre = Vec::new();
    for m in 0..p.vCoef.len() {
        let (ini, fin) = (p.vOff[m], p.vOff[m + 1]);
        vPre.clear();
        let mut acum = p.vCoef[m];
        for t in ini..fin {
            let x = vX[p.vIdVar[t]];
            let base = x.powi(p.vExp[t] - 1);
            vPre.push(acum * p.vExp[t] as f64 * base);
            acum *= base * x;
        }
        let mut suf = 1.0;
        for t in (ini..fin).rev() {
            vGrad[p.vIdVar[t]] += vPre[t - ini] * suf;
            suf *= vX[p.vIdVar[t]].powi(p.vExp[t]);
        }
    }
    vGrad
}
\end{lstlisting}
Rust exige \texttt{\#![allow(non\_snake\_case)]} para respetar tu nomenclatura camelCase, porque su convención es \texttt{snake\_case}.

\alt
\begin{itemize}
\item En Rust, la paralelización natural es \texttt{rayon} (\texttt{par\_iter}); no existe OpenMP nativo, y CUDA desde Rust depende de crates poco maduros \pv. Por eso la paralela se queda en C.
\item \texttt{nvcc} compila el código CUDA como C++. Tu ``CUDA en C'' será C++ con estilo C: \texttt{malloc}, punteros y sin clases. Es normal y aceptado.
\item Las comprobaciones de límites de Rust pueden costar algo de tiempo; los iteradores suelen eliminarlas. Mídelo y explícalo en la memoria en lugar de suponerlo.
\end{itemize}

\concepto{3.5 Opciones de compilación}
\que Las opciones que controlan optimización, avisos y diagnóstico.

\porque Las comparaciones de rendimiento solo son válidas con las mismas opciones en todas las versiones.

\como \texttt{-O3 -Wall -Wextra -std=c11} para medir, \texttt{-g -fsanitize=address,undefined} para los tests y \texttt{-fopenmp} para OpenMP. En CUDA: \texttt{nvcc -O3 -arch=sm\_XY -Xptxas=-v}. Todo el código de la guía compila sin avisos con \texttt{-Wall -Wextra}.

\alt \texttt{-march=native} activa AVX2 o AVX-512 y suele acelerar, pero el binario solo vale para la máquina donde compilas: compila en el servidor. \texttt{-ffast-math}: ver la trampa del concepto 1.4. Clang en lugar de gcc: válido para comparar.

\section{Rendimiento y medida}\label{tema:rend}
\soporte{F7, F8, F9 y F11}

\concepto{4.1 Medición de tiempos}
\que Tiempo de pared (\emph{wall-clock}, lo que tarda en terminar) frente a tiempo de CPU (suma de lo que consumen todos los hilos).

\porque El speedup se calcula con tiempo de pared. El enunciado aconseja \texttt{omp\_get\_wtime()} en CPU y pide desglosar transferencias y cómputo en GPU.

\como \texttt{omp\_get\_wtime()} en OpenMP, \texttt{clock\_gettime(CLOCK\_MONOTONIC)} en el secuencial de C, \texttt{std::chrono::steady\_clock} en C++, \texttt{std::time::Instant} en Rust y \texttt{cudaEventRecord}/\texttt{cudaEventElapsedTime} en CUDA (tema \ref{tema:cuda}). Mide solo la región de interés y reporta aparte lectura, parseo y cálculo.

\begin{trampa}
\texttt{clock()} devuelve tiempo de CPU acumulado de todos los hilos. Con 8 hilos marca unas 8 veces el tiempo real y \emph{oculta} el speedup. Tampoco midas un kernel CUDA con un reloj de CPU sin \texttt{cudaDeviceSynchronize()}: el lanzamiento es asíncrono.
\end{trampa}

\alt \texttt{gettimeofday} (obsoleto y no monótono), \texttt{rdtsc} (ciclos, poco portable) o \texttt{perf stat} (útil para contadores, no para tiempos por región).

\concepto{4.2 Speedup, eficiencia y escalabilidad}
\que Speedup $S_p=T_1/T_p$; eficiencia $E_p=S_p/p$. \emph{Escalado fuerte}: problema fijo y más hilos. \emph{Escalado débil}: problema proporcional a los hilos; se mide si $T_p$ se mantiene constante.

\porque El enunciado pide un \emph{análisis de escalabilidad} con \emph{análisis crítico}: no basta la tabla, hay que explicar por qué la eficiencia cae.

\como $T_1$ debe ser el \textbf{mejor secuencial}, no la versión paralela ejecutada con 1 hilo. Medida orientativa ($N=10^3$, $M=10^5$, $F=853\,478$, máquina de nube de 2 vCPU, gcc \texttt{-O2}, mediana de 3 repeticiones):

\begin{center}\small
\begin{tabular}{@{}lrrrr@{}}
\toprule
Variante & 1 hilo (s) & 2 hilos (s) & Speedup 2 hilos & Eficiencia \\
\midrule
Secuencial prefijo-sufijo & 0.0138 & -- & 1.00 & -- \\
\emph{reduction} (ingenuo $O(k^2)$) & 0.0801 & 0.0402 & 0.34 & 0.17 \\
\emph{atomic} (ingenuo $O(k^2)$) & 0.0806 & 0.0456 & 0.30 & 0.15 \\
Dos fases (prefijo-sufijo) & 0.0203 & 0.0109 & 1.27 & 0.63 \\
\bottomrule
\end{tabular}
\end{center}

Conclusiones que se repetirán en el servidor: un algoritmo peor no se arregla con más hilos, y la versión de dos fases con 1 hilo es más lenta que el secuencial porque hace una pasada de memoria extra (\texttt{vContrib}). Su beneficio aparece al crecer los hilos.

\alt Reportar \emph{throughput} (factores por segundo) en lugar de tiempo. Es útil para comparar tamaños distintos y lo recomiendo como columna adicional.

\concepto{4.3 Leyes de Amdahl y Gustafson}
\que Amdahl: si una fracción $s$ es secuencial, $S_p\le 1/\big(s+(1-s)/p\big)$, con techo $1/s$. Gustafson: si el problema crece con $p$, el speedup escalado es $S_p=s+(1-s)\,p$.

\porque Explican los resultados. Si el parseo (secuencial) es el 20\,\% del tiempo total, el speedup global nunca pasa de 5, aunque el cálculo escale perfectamente.

\como Mide la fracción secuencial real (lectura, transposición CSC, reserva de memoria) y superpón la curva de Amdahl en la gráfica de speedup.

\alt Ley de Karp-Flatt (fracción serie experimental $e=(1/S_p-1/p)/(1-1/p)$): si $e$ crece con $p$, el problema es la sobrecarga paralela y no la parte secuencial. Es muy útil para el análisis crítico.

\concepto{4.4 Jerarquía de memoria, intensidad aritmética y roofline}
\que La intensidad aritmética $I$ es el número de operaciones en coma flotante por byte movido desde memoria. El modelo roofline (Williams et al., 2009) acota el rendimiento: $P\le\min(P_{\text{pico}},\,I\cdot B)$, con $B$ el ancho de banda.

\porque Nuestro problema tiene $I$ baja: por factor se mueven unos 16 a 24 bytes (índice, exponente, lectura de \texttt{x} y, en dos fases, escritura de \texttt{vContrib}) para unas 10 operaciones, es decir, $I\approx0.4$ a $0.6$ flop/byte. Está \textbf{limitado por memoria} en CPU, y probablemente en GPU con \texttt{float}. Con \texttt{double} en una GeForce, el FP64 reducido puede acercar el límite al cómputo \pv. Este diagnóstico decide las optimizaciones: reducir bytes (tipos más pequeños, evitar la segunda pasada) antes que reducir operaciones.

\como Mide el ancho de banda con STREAM (\emph{triad}) en el servidor y compara con tus bytes por segundo efectivos ($F\cdot$bytes por factor$/T$). Si estás cerca del 60 a 80\,\% de STREAM, ya no hay mucho margen.

\alt Contar ciclos o instrucciones con \texttt{perf stat} (IPC, fallos de caché) da una visión complementaria a nivel de CPU.

\concepto{4.5 Cuellos de botella de este problema}
\begin{itemize}
\item \textbf{Accesos irregulares a \texttt{vX}}: \texttt{vIdVar} es aleatorio. Con $N=10^4$, \texttt{vX} ocupa 80 KB y cabe en L2; con $N=10^5$ ocupa 800 KB y puede no caber. Mide el tiempo por factor variando $N$ con $F$ fijo.
\item \textbf{Conflictos de escritura en \texttt{vGrad}}: varios monomios contribuyen a la misma variable. Con $N$ pequeño la contención de atómicos se dispara.
\item \textbf{Desequilibrio de carga}: $k_m$ varía de 1 a 16, así que \texttt{schedule(static)} puede repartir mal.
\item \textbf{Núcleos heterogéneos}: el i5-13420H mezcla núcleos P y E \pv. Con reparto estático, los hilos en núcleos E terminan después y frenan a todos.
\item \textbf{\emph{False sharing}}: arrays privados por hilo contiguos que comparten línea de caché (64 bytes).
\item \textbf{PCIe}: subir 80 MB a la GPU a unos 10 a 25 GB/s cuesta de 3 a 8 ms \pv, probablemente más que el kernel. Por eso hay que desglosar los tiempos.
\end{itemize}

\concepto{4.6 Perfiladores}
\que Herramientas que miden dónde se va el tiempo.

\porque Para justificar optimizaciones con datos, como pide la parte opcional de CUDA (\emph{análisis profundo de la ejecución}).

\como CPU: \texttt{perf stat -e cycles,instructions,cache-misses ./omp} y \texttt{perf record}/\texttt{perf report}. GPU: \texttt{nsys profile ./grad} para la línea temporal (transferencias frente a kernels) y \texttt{ncu --set full ./grad} para métricas por kernel (ocupación lograda, ancho de banda, divergencia).

\alt \texttt{gprof} (antiguo, instrumenta), Valgrind \texttt{cachegrind} (muy lento, simula cachés) e Intel VTune.

\begin{trampa}
\texttt{nvprof} y Visual Profiler son obsoletos y no admiten GPUs de capacidad 8.0 o superior. Usa Nsight Systems y Nsight Compute.
\end{trampa}
\clearpage
